\input{myquizpreamble}
\input{yliow}
\input{ciss362}
\textwidth=6in

\renewcommand\TITLE{Assignment a11}


\renewcommand\AUTHOR{John Doe}

\begin{document}
\topmatter

\textsc{Objectives}
\begin{itemize}
\li Design TM
\end{itemize}

Here are the questions:
\begin{myenum}
\li Sipser 3.1. Q1. Solution to 3.1(b) is provided in the textbook. Study it carefully. 
\li Sipser 3.2. Study the solution to 3.2(a).
\li Sipser 3.3. Study the solution in the textbook.
\li Sipser 3.5. Study the solution in the textbook.
\li Sipser 3.7. Q2.
\li Sipser 3.8. Q3. Solution to 3.8(a) is provided in the textook. Study it carefully.
\li Sipser 3.9. Q4. With modification. 
\li Sipser 3.10. Solution to 3.10 is provided in the textook. Study it carefully. 
\li Sipser 3.11. Q5. 
\li Sipser 3.15. Q6. Solution to 3.15(a) is in the textbook. Study it carefully.
\li Sipser 3.16. Q7. Solution to 3.16(a) is in the textbook. Study it carefully.
Ignore (e).
\end{myenum}


%------------------------------------------------------------------------------
\newpage
\textsc{How to draw a state diagram}

Here's an example showing you how to draw the elements of a state diagram.
Also, look at the solution to 1.3 below.

\begin{center}
  \begin{tikzpicture}[shorten >=1pt,>=triangle 60,double distance=2pt,node distance=2cm,auto,initial text=]
    
    \node[state,initial]   (q0) at (0, 0) {$q_0$};
    \node[state]           (q1) at (4, 0) {$q_1$};
    \node[state,accepting] (q2) at (8, 0) {$q_2$};
    \node[state]           (q3) at (0,-4) {$q_3$};
    \node[state]           (q4) at (4,-4) {$q_4$};

    \node[state]           (q5) at (12, 0) {$q_5$};
    \node[state]           (q6) at (12,-4) {$q_6$};

    \path[->]
    (q0) edge                node {$1$} (q1)    
    (q1) edge [bend left=15] node {$0$} (q2)
    (q0) edge [left]  node {$\alpha$} (q3)
    (q1) edge [right] node {$\beta$}  (q4)
    (q1) edge [loop above]   node {$1$} (q1)
    (q2) edge [loop below]   node {$0$} ()
    (q2) edge [bend left=15] node {$1$} (q1)
    (q3) edge [loop right]   node {$0,1$} ()
    (q4) edge [loop left]    node {$2,3$} ()
    (q5) edge [right, bend left=30]    node {$\gamma$} (q6)
    (q5) edge [right, bend left=60]    node {$\ep$}    (q6)
    (q5) edge [left, bend right=30]    node {$\delta$} (q6)
    ;
  \end{tikzpicture}
\end{center}

For more information on drawing state diagrams go to my tutorials and
look for \verb!latex-automata.pdf!:
\begin{center}
{\small \url{https://drive.google.com/file/d/1AeE-P0WNvQlitzPDxQpGE8bMR9Yc9gMW}}
\end{center}
Let me know if you have any questions about drawing state diagram.

%------------------------------------------------------------------------------
\textsc{How to write the computation done by a TM}

A TM computes just like a DFA/NFA/PDA computes.
A snapshot of the state of computation of a TM is called (just like the
case for DFA) a \lq\lq instantaneous description" or
\lq\lq configuration" for that TM.
For instance is I say the following
\[
Aq_1aabbb
\]
is the \lq\lq configuration" or \lq\lq instantaneous description"
of some TM while doing a computation, I mean
at this point in the computation,
\begin{itemize}
\li the TM has the string $Aaabbb$ (followed by an
infinitely number of spaces) in it's input tape
\li the TM is at state $q_1$
\li the TM has its read/write head at the leftmost $a$ (i.e., index $1$).
The character that the read/write head is about to read is always
the character on the right of the state in the configuration:
\[
Aq_1\underline{a}abbb
\]
\end{itemize}

Here's how to write a derivation of \lq\lq instantaneous descriptions" or
\lq\lq configurations" for a TM:
\begin{align*}
q_0 aaabbb &\vdash Aq_1aabbb \\ 
           &\vdash Aaq_1abbb \\ 
           &\vdash Aaaq_1bbb \\
           &\vdash AaaBq_2bb 
\end{align*}
where the following is an incomplete list of transitions for this TM:
\begin{align*}
\delta(q_0, a) &= (q_1, A, R) \\
\delta(q_1, a) &= (q_1, a, R) \\
\delta(q_1, b) &= (q_2, B, R) 
\end{align*}

%------------------------------------------------------------------------------
\newpage
Q1. Sipser 3.1. 
Solution to 3.1(b) is provided in the textbook. Study it carefully.
Write the TM computations using the \lq\lq $\vdash$" notation.
\begin{myenum}
\li Do \verb!$\vdash$! to get $\vdash$.
\li Do \verb!$\SPACE$! to get $\SPACE$.
\li Do \verb!\texttt{x}! to write x in typewriter font: \texttt{x}.
\li Do \verb!$q_{\text{accept}}$! to get $q_{\text{accept}}$.
\li Do \verb!$q_{\text{reject}}$! to get $q_{\text{reject}}$.
\end{myenum}

\textsc{Solution}.


(a)

(b)
\begin{align*}
  q_1\texttt{00} &\vdash \SPACE q_2 \texttt{0} \\
        &\vdash \SPACE \texttt{x} q_3 \SPACE \\
        &\vdash \SPACE q_5 \texttt{x} \SPACE \\
        &\vdash q_5 \SPACE \texttt{x} \SPACE \\
        &\vdash \SPACE q_2 \texttt{x} \SPACE \\
        &\vdash \SPACE \texttt{x} q_2 \SPACE \\
        &\vdash \SPACE \texttt{x} \SPACE q_{\text{accept}} \\
\end{align*}

(c)

(d)


%------------------------------------------------------------------------------
\newpage
Q2. Sipser 3.7.

\textsc{Solution}.

\input{q02.tex}

%------------------------------------------------------------------------------
\newpage
Q3. Sipser 3.8.
Solution to 3.8(a) is provided in the textook. Study it carefully.

For each part, design a TM and then
implement your TMs using the python
program I gave, test it fully, and then paste your TM program into the
solution area of this question.

(d)
Recall that
the language
\[
L = \{w \mid w \in \{a, b, c\}^*, \,\,\,
\text{$w$ has the same number of $a$'s, $b$'s, $c$'s} \}
\]
is not a CFL.
Design a TM that accepts $L$.

(e) \textsc{Optional DIY:}
Design a TM that checks if an integer is prime.
To check is $n$ is a prime, the user enters
$\langle n \rangle$ (which is $n$ 1s) on the input tape.
What about checking that the input is composite?
(A composite number is an integer $> 1$ that is not a prime.)

\textsc{Solution}.


(b)
\begin{answercode}

\end{answercode}
(Put your code above.)


(c)
\begin{answercode}

\end{answercode}


(d)
\begin{answercode}

\end{answercode}

(e)
\begin{answercode}

\end{answercode}



%------------------------------------------------------------------------------
\newpage
Q4. Sipser 3.9.

For (a), show that a 2-PDAs is more powerful than 1-PDAs and in fact
2-PDAs are Turing machines.

\textsc{Solution}.


(a)

(b)


%------------------------------------------------------------------------------
\newpage
Q5. Sipser 3.11.

(Basically, you need to simulate the left-hand-side tape on the right-hand-side tape.
Right?
You can take a pair of scissors and snip the two-way infinite tape into two pieces
and then try to join the left-hand-side tape onto the right-hand-side tape.
But you cannot stick the left-hand-side infinte tape onto the right-hand-side \textit{infinite} tape.
Because in that case you have to travel infinitely far to the right
to reach the left-hand-side tape, which means you can never reach it.
So what should you do?)

\textsc{Solution}.

\input{q05.tex}


%------------------------------------------------------------------------------
\newpage
Q6. Sipser 3.15. 
Solution to 3.15(a) is in the textbook. Study it carefully.

\textsc{Solution}.


(b)

(c)

(d)

(e)


%------------------------------------------------------------------------------
\newpage
Q7.
Solution to 3.16(a) is in the textbook. Study it carefully.
Ignore (e).

\textsc{Solution}.


(b)

(c)

(d)


\end{document}
